\section{Rúbrica de Evaluación y Autoevaluación}

\begin{table}[h!]
\centering
\small
\renewcommand{\arraystretch}{1.3}
\begin{tabular}{|p{4.5cm}|c|p{7cm}|}
\hline
\textbf{Criterio de Evaluación} & \textbf{Nota} & \textbf{Justificación Técnica y Detalle} \\ \hline
\textbf{1. Funcionalidad general} \newline (Navegación, persistencia, UX) & Bueno (5) & Todos los flujos principales (autenticación JWT, listado de cartelera, compra asíncrona, cancelación y transferencia) están operativos, conectados a un VPS activo y con refresco inmediato. \\ \hline
\textbf{2. Interfaz de usuario} \newline (Diseño, maquetación, layouts) & Bueno (5) & Diseño oscuro premium con acentos vibrantes, widget de boleto troquelado, generación dinámica de código QR y adaptabilidad completa. \\ \hline
\textbf{3. Documentación} \newline (Documento LaTeX + video) & Bueno (5) & Presentación modular en LaTeX, documentación detallada paso a paso y enlace al video demostrativo de ejecución. \\ \hline
\textbf{4. Código y buenas prácticas} \newline (Modularidad, commits, git) & Bueno (5) & Patrón arquitectónico MVVM, métodos auxiliares limpios fuera del build, commits semánticos ordenados y README estructurado. \\ \hline
\textbf{Puntaje Autoevaluado Total} & \multicolumn{2}{c|}{\textbf{20 / 20 puntos}} \\ \hline
\end{tabular}
\caption{Tabla de Autoevaluación del Proyecto Final.}
\label{tab:rubrica}
\end{table}
